\documentclass[english, parskip=full]{scrartcl}
\usepackage[utf8]{inputenc}  % Kodierung der Datei
\usepackage[english]{babel}  % Sprache auf Deutsch 
\usepackage{color}
\usepackage{graphicx, gensymb}
\usepackage{amsfonts, cancel, amssymb}
\usepackage{amsmath, amsthm, array, esint}
\usepackage[most]{tcolorbox}
\usepackage{booktabs, multirow}  % schönere Tabellen
\usepackage{caption}  % Anpassung der Captions
\usepackage{scrlayer-scrpage}    % Manipulation des Seitenstils
\pagestyle{scrheadings}  % Stil für die Seite setzen
\automark{section}  % Abschnittsnamen als Seitenbeschriftung 
\setheadsepline{.5pt}    % Kopzeile durch Linie abtrennen
\usepackage[hidelinks]{hyperref}  % Links und weitere PDF-Features
\usepackage{typearea, tikz, pgffor, verbatim}
\usetikzlibrary{calc, arrows.meta,arrows, graphs, quotes, positioning}
\usepackage[cache=false, outputdir=build]{minted}
\usemintedstyle{vs}
\usepackage{placeins} % provides \FloatBarrier

\definecolor{bg}{rgb}{0.9,0.9,0.9}
\newcommand\numberthis{\addtocounter{equation}{1}\tag{\theequation}}
\usepackage{svg}
\usepackage{siunitx}
\usepackage{physics}
\usepackage{tensor}
\usepackage{slashed}
\usepackage{bbm}
\renewcommand{\bra}{}
\newcommand{\kom}{}
\newcommand{\brak}{}
\renewcommand{\braket}{}
\newcommand{\intx}{}
\newcommand{\intr}{}
\newcommand{\kla}{}
\newcommand{\klam}{}
\newcommand{\klammer}{}
\newcommand{\oper}{}
\newcommand{\tecxt}{}
\newcommand{\Bra}{}
\newcommand{\Ket}{}
\renewcommand{\i}{\text{i}}
\newcommand{\e}{\text{e}}
\newcommand*{\Fourier}{\textsc{Fourier}\ }
\newcommand*{\F}{\mathcal{F}}
\newcommand*{\sinc}{\text{sinc\,}}
\newcommand{\conv}{\mathop{\scalebox{3.5}{\raisebox{-0.38ex}{\makebox[0.5\width][c]{$\ast$}}}}\limits}

\newcommand*{\titel}{06 Quantization of Free Fields}
\newcommand*{\autor}{Joachim Schwardt}

\hypersetup{pdfauthor={\autor}, pdftitle={\titel}}

%\titlehead{titlehead}
\subject{Relativistic Quantum Field Theory}
\title{\titel}
\author{\autor}
\date{\today}

\begin{document}

%\begin{titlepage}
\maketitle  % Titel setzen
%\tableofcontents  % Inhaltsverzeichnis setzen
%\end{titlepage}

\section{Short exercise on propagators}
The time-ordered product of two fermionic fields is given by 
\begin{align}
\mathcal{T} \hat{\psi}(x) \bar{\hat{\psi}}(y) := \Theta\kla\left( x^0 - y^0 \right) \hat{\psi}(x) \bar{\hat{\psi}}(y) - \Theta \kla\left( y^0 - x^0 \right) \bar{\hat{\psi}}(y) \hat{\psi}(x).
\end{align}
%Note that for $\partial_\mu \equiv \frac{\partial}{\partial x^\mu}$ we find
%\begin{align}
%\slashed{\partial} \bar{\hat{\psi}}(y) &= \gamma^\mu \hat{\psi}^\dagger (y) \gamma^0 \partial_\mu = \hat{\psi}^\dagger (y) \kla\left( \klammer\left\{ \gamma^\mu, \gamma^0 \right\} - \gamma^0 \gamma^\mu  \right) \partial_\mu = 2\bar{\hat{\psi}}(y) \gamma^0 \partial_0 -\bar{\hat{\psi}} (y) \slashed{\partial} .
%\end{align}
%Then we find, using the equation of motion for the fields $\left( \i\slashed{\partial} - m \right) \hat{\psi}(x) = 0$
%\begin{align*}
%\kla\left( \i\slashed{\partial} - m \right) \braket\langle 0 | \mathcal{T} \hat{\psi}(x) \bar{\hat{\psi}}(y) | 0 \rangle &= \Theta\kla\left( x^0 - y^0 \right) \braket\langle 0 | \klam\left[ \left( \i\slashed{\partial} - m \right) \hat{\psi}(x) \right] \bar{\hat{\psi}}(y) | 0 \rangle - \\
%&\quad - \Theta\kla\left( y^0 - x^0 \right) \braket\langle 0 | \bar{\hat{\psi}}(y) \kla\left( 2\i\gamma^0\partial_0 - \i \slashed{\partial} - m \right) \hat{\psi}(x) | 0 \rangle = \\
%&= 0 - 2\Theta \kla\left( y^0 - x^0 \right) \braket\langle 0 | \bar{\hat{\psi}} (y) \left( \i\gamma^0\partial_0 - m \right) \hat{\psi}(x) | 0 \rangle = 0.
%\end{align*}
Note that
\begin{align*}
\partial_0 \langle 0 | \mathcal{T} \hat{\psi}(x) \bar{\hat{\psi}}(y) | 0 \rangle &= \delta\kla\left( x^0 - y^0 \right) \langle 0 | \hat{\psi}(x) \bar{\hat{\psi}}(y) | 0 \rangle + \Theta\kla\left( x^0 - y^0 \right) \langle 0 | \partial_0 \hat{\psi}(x) \bar{\hat{\psi}}(y) | 0 \rangle +\\
&\quad + \delta \kla\left( x^0 - y^0 \right) \langle 0 | \bar{\hat{\psi}}(y) \hat{\psi}(x) | 0 \rangle - \Theta\kla\left( y^0 - x^0 \right) \langle 0 | \bar{\hat{\psi}}(y) \partial_0 \hat{\psi}(x) | 0 \rangle = \numberthis\\
&= \delta\kla\left( x^0 - y^0 \right) \braket\langle 0 | \klammer\left\{ \hat{\psi}(x), \bar{\hat{\psi}}(y) \right\} | 0 \rangle + \langle 0 | \mathcal{T} \partial_0 \hat{\psi}(x) \bar{\hat{\psi}}(y) | 0 \rangle  = \numberthis \\
&= \delta (x - y) + \langle 0 | \mathcal{T} \partial_0\hat{\psi}(x) \bar{\hat{\psi}}(y) | 0 \rangle. \numberthis
\end{align*}
Thus we find
\begin{align}
\slashed{\partial} \langle 0 | \mathcal{T} \hat{\psi}(x) \bar{\hat{\psi}}(y) | 0 \rangle &= \delta (x-y)\gamma^0 + \langle 0 | \mathcal{T} \slashed{\partial} \hat{\psi}(x) \bar{\hat{\psi}}(y) | 0 \rangle
\end{align}
and finally using the equation of motion
\begin{align}
\kla\left( \i\slashed{\partial} - m \right) \langle 0 | \mathcal{T} \hat{\psi}(x) \bar{\hat{\psi}}(y) | 0 \rangle &= \delta (x-y)\gamma^0 + \langle 0 | \mathcal{T} \left( \i\slashed{\partial} - m \right) \hat{\psi}(x) \bar{\hat{\psi}}(y) | 0 \rangle = \delta (x-y)\gamma^0.
\end{align}
\section{Free fermionic spin-$\frac{1}{2}$ field in the Kugo-Ojima formalism}
We want to consider the free QED theory with the inclusion of Faddeev-Popov ghosts.
We will quantize the free spin-$\frac{1}{2}$ electron field, the free photon field and the free ghost fields.
Gauge invariance will be replaced by BRST invariance, which remains valid after gauge fixing.
The Noether theorem will lead to a conserved BRST charge described by an operator $\hat{Q}_B$.
The Gupta-Bleuler formalism will be replaced by the Kugo-Ojima formalism which is more general and applicable to non-abelian gauge theories.
The entire setup appears in full QED as the description of asymptotically free particles long before or after interactions.\\
The Lagrangian is given by
\begin{align}
\mathcal{L} &= \mathcal{L}_\tecxt\text{g.inv.} + \mathcal{L}_\tecxt\text{fix} + \mathcal{L}_\tecxt\text{ghost},\\
\mathcal{L}_\tecxt\text{g.inv.} &= -\frac{1}{4} F^{\mu\nu} F_{\mu\nu} + \i \bar{\psi} \slashed{\partial} \psi,\\
\mathcal{L}_\tecxt\text{fix+ghost} &= s\klam\left[ \bar{c} \kla\left( \partial_\mu A^\mu  \right) + \frac{\xi}{2}B \right] = B (\partial_\mu A^\mu ) + \frac{\xi}{2}B^2 - \bar{c} \square c.
\end{align}
Here $B(x)$ is a bosonic field and $c(x),\bar{c}(x)$ are fermionic ghost and anti-ghost fields respectively.
All are treated as scalar fields under Poincar\'e transformations.
On the classical level $c$ and $\bar{c}$ will be treated as anti-commutation number valued functions.
On the quantum level we will instead require canonical anti-commutation relations similar to those for the electro field.
The BRST transformations generalize gauge transformations and in the free theory may be expressed using the operator $s$ 
\begin{align}
sA^\mu (x) &= \partial^\mu c(x),\\
s\psi_j(x) &= 0,\\
s \bar{\psi}_j(x) &= 0,\\
sc(x) &= 0,\\
s\bar{c}(x) &= B(x),\\
sB(x) &= 0
\end{align}
where $s$ should be interpreted as a fermionic differential operator satisfying product rules
\begin{align}
s \kla\left( B_1B_2 \right) &= \kla\left( sB_1 \right)B_2 + B_1 \kla\left( sB_2 \right),\\
s \kla\left( F_1B_2 \right) &= \kla\left( sF_1 \right) B_2 - F_1 \kla\left( sB_2 \right),\\
s \kla\left( F_1F_2 \right) &= \kla\left( sF_1 \right) F_2 - F_1 \kla\left( sF_2 \right).
\end{align}
We will quantize this theory for the special case of $\xi=1$, determine the conserved charge $\hat{Q}_B$, evaluate the Fock space of states on which all operators are defined, determine the 1-particle states of positive and negative or indefinite norm, and compute the commutation relations of $\hat{Q}_B$ with all creation and annihilation operators.
\subsection{Constraints of the theory}
There are canonical momenta for $A^\mu$, $B$ and $c$.
Using $F^{\mu\nu} = \partial^\mu A^\nu - \partial^\nu A^\mu$ we find
\begin{align}
\pi_\rho &= \frac{\partial \mathcal{L}}{\partial \dot{A}^\rho} = \\
&= -\frac{1}{4} \frac{\partial}{\partial \dot{A}^\rho} \klam\left[ 2  \kla\left( \partial^\mu A^\nu \right) \partial_\mu A_\nu - 2 \kla\left( \partial^\nu A^\mu  \right) \partial_\mu A_\nu \right] + B \frac{\partial \dot{A}^0}{\partial \dot{A}^\rho} = \\
&= -\frac{1}{2} \klam\left[ 2\delta_{0}^{\ \mu} \delta_{\rho}^{\ \nu} \partial_\mu A_\nu - 2\delta_{0}^{\ \nu}\delta_{\rho}^{\ \mu} \partial_\mu A_\nu \right] + B\delta_{0\rho} = \\
&= - \klam\left[ \partial_0 A_\rho - \partial_\rho A_0 \right] + B\delta_{0\rho} = \\
&= \begin{cases} B &\tecxt\text{for } \rho = 0,\\
F_{\rho 0} &\tecxt\text{for } \rho = 1,2,3 \end{cases},\\
\pi_B &= \frac{\partial \mathcal{L}}{\partial \dot{B}} = 0,\\
\pi_c &= \frac{\partial \mathcal{L}}{\partial \dot{c}} = \dot{\bar{c}},\\
\pi_{\bar{c}} &= \frac{\partial \mathcal{L}}{\partial \dot{\bar{c}}} = \dot{c}.
\end{align}
%Since $\pi_0 = B$ has no dependence of the dynamic variable $\dot{A}^\mu$, this condition represents a constraint. 
\subsection{Eliminating $B$}
The equation of motion for $B$ gives
\begin{align}
0 &= \partial_\mu \frac{\partial \mathcal{L}}{\partial \partial_\mu B} - \frac{\partial \mathcal{L}}{\partial B} = 0 - \partial_\mu A^\mu - B
\end{align}
giving $B = -\partial_\mu A^\mu$.
Since the Lagrangian does not contain any dynamic $B$-part, the constraint is of second type and $B$ may be treated merely as an abbreviation.
Then the Lagrangian becomes
\begin{align}
\mathcal{L} &= -\frac{1}{4} F^{\mu\nu} F_{\mu\nu} + \i \bar{\psi} \slashed{\partial}  \psi - \frac{1}{2} \kla\left( \partial_\mu A^\mu  \right)^2 - \bar{c} \square c.
\end{align}
The non-trivial changes to the BRST transformations are
\begin{align}
s\bar{c}(x) &= - \partial_\mu A^\mu,\\
s\partial_\mu A^\mu(x) &= \partial_\mu sA^\mu(x) = -\square c(x).
\end{align}
\subsection{Noether current for BRST invariance}
We consider the infinitesimal field variation $\delta\phi = \theta s\phi$ for all fields $\phi$ where $\theta$ is an infinitesimal fermionic number.
Since this variation is zero for $\phi \in \klammer\left\{ \psi, \bar{\psi}, c \right\}$, the only non-trivial transformations are $\phi \in \klammer\left\{ A_\mu, \bar{c} \right\}$.
Let us first consider $\phi = \bar{c}$.
Then the relevant part of the Lagrangian is
\begin{align}
\mathcal{L}' &= - \kla\left( \bar{c} + \theta s\bar{c} \right) \square c = \mathcal{L} + \theta \partial_\mu A^\mu \square c.
\end{align}
On its own this part of the Lagrangian is therefore not invariant, but it need not be.
For $\phi = A_\mu$ we find
\begin{align}
\mathcal{L}' &= -\frac{1}{2} \klam\left[  \kla\left( \partial^\mu \kla\left( A^\nu + \theta sA^\nu \right) - (\mu\leftrightarrow\nu) \right) \partial_\mu \kla\left( A_\nu + \theta s A_\nu \right) \right] - \frac{1}{2} \kla\left( \partial_\mu A^\mu + \partial_\mu \theta sA^\mu  \right)^2 = \\
&= \mathcal{L} - \frac{1}{2} \klam\left[ \kla\left( \partial^\mu \theta \partial^\nu c - \partial^\nu \theta \partial^\mu c  \right) \partial_\mu A_\nu + \kla\left( \partial^\mu A^\nu - \partial^\nu A^\mu  \right) \partial_\mu \theta \partial_\nu c \right] - \partial_\mu A^\mu \partial_\nu \theta \partial^\nu c = \\
&= \mathcal{L} - \theta \klammer\left\{ \kla\left( \partial^\mu \partial^\nu c \right) \left(\partial_\mu A_\nu - \partial_\nu A_\mu \right) + \partial_\mu A^\mu \square c \right\} = \\
&= \mathcal{L} - \theta \partial_\mu A^\mu \square c.
\end{align}
This is precisely the opposite change as compared to the $\phi = \bar{c}$-case.
Thus the Lagrangian is indeed invariant under the full BRST transformation.
The Noether current is given by
\begin{align}
j^\rho &= \sum_\phi \frac{\partial \mathcal{L}}{\partial \partial_\rho \phi} \delta \phi = \\
&= \frac{\partial \mathcal{L}}{\partial \partial_\rho \bar{c}} \theta s\bar{c} + \frac{\partial \mathcal{L}}{\partial \partial_\rho A_\mu} \theta sA_\mu = \\
&= \partial^\rho c \theta \kla\left( -\partial_\mu A^\mu  \right) - \klam\left[ \delta^{\rho \nu} \delta^{\mu \lambda} \partial_\nu A_\lambda - \delta^{\rho \lambda} \delta^{\mu \nu} \partial_\nu A_\lambda \right] \theta \partial_\mu c - \partial_\nu A^\nu \delta^{\rho}_{\ \lambda} \delta^{\mu \lambda} \theta \partial_\mu c = \\
&= +\theta \partial_\mu A^\mu \partial^\rho c - \klam\left[ \partial^\rho A^\mu - \partial^\mu A^\rho  \right] \theta \partial_\mu c - \partial_\nu A^\nu \theta \partial^\rho c  = \\
&= -\theta F^{\rho \mu} \partial_\mu c.
\end{align}
The conserved charge corresponding to the symmetry is
\begin{align}
Q_B &:= \intr\int_{\mathbb{R}^{3}} j^0 \, \mathrm{d}^{3}x = -\theta \intr\int_{\mathbb{R}^{3}} F^{0j} \partial_j c \, \mathrm{d}^{3}x = \theta \intr\int_{\mathbb{R}^{3}} \kla\left( \partial_\mu F^{0\mu} \right) c \, \mathrm{d}^{3}x.
\end{align}
Note that the equation of motion for the $A_\mu$ is given by
\begin{align}
0 &= \partial_\rho \frac{\partial \mathcal{L}}{\partial \partial_\rho A_\mu} - \frac{\partial \mathcal{L}}{\partial A_\mu} = \\
&= \partial_\rho \klam\left[ \partial^\rho A^\mu - \partial^\mu A^\rho - \partial_\nu A^\nu \delta^{\mu\rho} \right] - 0 = \partial_\rho F^{\rho \mu} - \partial^\mu \partial_\nu A^\nu.
\end{align}
Thus we have $\partial_\mu F^{0\mu} = -\partial^0 \partial_\nu A^\nu$ and 
\begin{align}
Q_B &= \theta \intr\int_{\mathbb{R}^{3}} \kla\left( -\partial^0 \partial_\nu A^\nu \right) c \, \mathrm{d}^{3}x \equiv \theta \intr\int_{\mathbb{R}^{3}} \dot{B}c \, \mathrm{d}^{3}x.
\end{align}


%\begin{thebibliography}{9}
%\bibitem{example} description, \url{https://www.exmapleurl}, day.\,month~2021.
%\end{thebibliography}

\end{document}